% TODO make hw stylesheet
\documentclass[10pt]{article}
\usepackage[letterpaper,top=1in]{geometry}
\usepackage[dvipsnames,svgnames]{xcolor}
\usepackage[shortlabels]{enumitem}
\usepackage[framemethod=TikZ]{mdframed}
\usepackage{amsmath,amssymb,amsthm}
\usepackage[colorlinks]{hyperref}
\usepackage{multicol}
\usepackage{tabularx}

\hypersetup{
	colorlinks=true,
	linkcolor=blue,
	filecolor=magenta,
	urlcolor=magenta,
	pdftitle={MAT 528 HW}
}

\usepackage{mathtools}
\usepackage{thmtools}
\usepackage[textsize=footnotesize]{todonotes}
\usepackage{titling}

\reversemarginpar
\mdfdefinestyle{mdbluebox}{roundcorner=10pt,innerbottommargin=9pt,
	linecolor=blue,backgroundcolor=TealBlue!5,}
\declaretheoremstyle[headfont=\sffamily\bfseries\color{MidnightBlue},
mdframed={style=mdbluebox},]{thmbluebox}

\mdfdefinestyle{mdbloobox}{frametitlefont=\bfseries,innerbottommargin=8pt,
	nobreak=true,backgroundcolor=TealBlue!5,linecolor=blue,}
\declaretheoremstyle[headfont=\bfseries\color{MidnightBlue},
mdframed={style=mdbloobox},headpunct={\\[3pt]},postheadspace=0pt,]{thmbloobox}

\mdfdefinestyle{mdredbox}{frametitlefont=\bfseries,innerbottommargin=8pt,
	nobreak=true,backgroundcolor=Salmon!5,linecolor=RawSienna,}
\declaretheoremstyle[headfont=\bfseries\color{RawSienna},
mdframed={style=mdredbox},headpunct={\\[3pt]},postheadspace=0pt,]{thmredbox}

\mdfdefinestyle{mdgreenbox}{linecolor=ForestGreen,backgroundcolor=ForestGreen!5,
	linewidth=2pt,rightline=false,leftline=true,topline=false,bottomline=false,}
\declaretheoremstyle[headfont=\bfseries\sffamily\color{ForestGreen!70!black},
mdframed={style=mdgreenbox},headpunct={ --- },]{thmgreenbox}

\mdfdefinestyle{mdorangebox}{linecolor=RedOrange,backgroundcolor=RedOrange!5,
	linewidth=2pt,rightline=false,leftline=true,topline=false,bottomline=false,}
\declaretheoremstyle[headfont=\bfseries\sffamily\color{RedOrange!70!black},
mdframed={style=mdorangebox},headpunct={ --- },]{thmorangebox}

\mdfdefinestyle{mdblackbox}{linecolor=black,backgroundcolor=RedViolet!5!gray!5,
	linewidth=3pt,rightline=false,leftline=true,topline=false,bottomline=false,}
\declaretheoremstyle[mdframed={style=mdblackbox}]{thmblackbox}

\declaretheoremstyle[spaceabove=6pt,spacebelow=6pt]{basehead}

\declaretheorem[style=basehead,name=Problem]{problem}
\declaretheorem[style=thmredbox,name=Problem,numbered=no]{problem*}
\declaretheorem[style=thmbloobox,name=Setup]{setup} % for config geo
\declaretheorem[style=thmredbox,name=Example,sibling=problem]{example}
\declaretheorem[style=thmredbox,name=$\bigstar$ Problem,sibling=problem]{niceprob}
\declaretheorem[style=thmbluebox,name=Theorem,sibling=problem]{theorem}
\declaretheorem[style=thmbluebox,name=Corollary,sibling=theorem]{corollary}
\declaretheorem[style=thmbluebox,name=Proposition,sibling=theorem]{prop}
\declaretheorem[style=thmbluebox,name=Lemma,numberwithin=problem]{lemma}
\declaretheorem[style=thmbluebox,name=Theorem,numbered=no]{theorem*}
\declaretheorem[style=thmbluebox,name=Lemma,numbered=no]{lemma*}
\declaretheorem[style=thmgreenbox,name=Claim,numberwithin=problem]{claim}
\declaretheorem[style=thmgreenbox,name=Claim,numbered=no]{claim*}
\declaretheorem[style=thmblackbox,name=Remark,numberwithin=problem]{remark}
\declaretheorem[style=thmblackbox,name=Remark,numbered=no]{remark*}
\declaretheorem[style=thmgreenbox,name=Definition,sibling=theorem]{definition}
\declaretheorem[style=thmgreenbox,name=Definition,numbered=no]{definition*}

\providecommand{\ol}{\overline}
\providecommand{\ceil}[1]{\left\lceil #1 \right\rceil}
\providecommand{\floor}[1]{\left\lfloor #1 \right\rfloor}
\newcommand{\dd}{\mathrm{d}}
\providecommand{\ul}{\underline}
\providecommand{\eps}{\varepsilon}
\providecommand{\half}{\frac{1}{2}}
\providecommand{\CC}{\mathbb C}
\providecommand{\FF}{\mathbb F}
\providecommand{\NN}{\mathbb N}
\providecommand{\QQ}{\mathbb Q}
\providecommand{\RR}{\mathbb R}
\providecommand{\ZZ}{\mathbb Z}
\providecommand{\dg}{^\circ}
\providecommand{\ii}{\item}
\providecommand{\setdiff}{\smallsetminus}
\providecommand{\alert}[1]{{\sffamily\textbf{\textcolor{blue}{#1}}}}
\providecommand{\inv}{^{-1}}
\providecommand{\mb}[1]{\mathbf{#1}}

\newenvironment{soln}{\begin{proof}[Solution]}{\end{proof}}

\providecommand{\pow}{\operatorname{Pow}}
\providecommand{\vocab}[1]{{\textbf{\textcolor{ForestGreen}{#1}}}}
\providecommand{\lcm}{\operatorname{lcm}}
\providecommand{\abs}[1]{\left\lvert #1\right\rvert}
\providecommand{\norm}[1]{\left\|#1\right\|}

\usepackage{mathrsfs}

% place title at top right
\setlength{\droptitle}{-8ex}
\pretitle{\begin{flushright}\Large\bfseries}
\posttitle{\par\end{flushright}}
\preauthor{\begin{flushright}}
\postauthor{\end{flushright}}
\predate{\begin{flushright}}
\postdate{\end{flushright}}

\title{MAT 528 HW10}
\author{\large Aditya Pahuja}
\date{\large Due 11/20}
\begin{document}
\maketitle

\begin{problem}
    [Munkres 51.2]
    Given spaces $X$ and $Y$, let $[X,Y]$ denote the set of
    homotopy classes of maps of $X$ into $Y$.
    \begin{enumerate}[(a)]
	\ii Let $I = [0,1]$. Show that for any $X$,
	the set $[X,I]$ has a single element.
	\ii Show that if $Y$ is path connected,
	the set $[I,Y]$ has a single element.
    \end{enumerate}
\end{problem}

\begin{soln}
    (a) Each continuous map $f\colon X\to I$ is homotopic to the zero map
    via the homotopy
    \begin{align*}
	F\colon X\times I&\to I\\
	F(x,t) &= (1-t)f(x).
    \end{align*}
    We can see that $F(x,0) = (1-0)f(x) = f(x)$ and $F(x,1) = (1-1)f(x) = 0$
    for all $x\in X$, so to verify that $F$ is a homotopy
    it remains to check continuity of $F$;
    this follows from the continuity of the maps
    $(x,t)\mapsto f(x)$ and $(x,t)\mapsto 1-t$
    and Munkres's Theorem 21.5.
    \\

    (b) Let $f,g\colon I\to Y$ be two paths in $Y$.
    Let $f(0) = x_0$, $f(1) = x_1$, $g(0) = y_0$, and $g(1) = y_1$.
    Since $Y$ is path connected, there exists a path
    $h\colon I\to Y$ from $x_1$ to $y_0$.
    Define $F\colon I\times I\to Y$ by
    \begin{equation*}
        F(x,t) = \begin{cases}
	    f(x + 2t) & \text{if }x+2t \le 1\\
	    h(x + 2t - 1) & \text{if }1\le x+2t \le 2\\
	    g(x + 2t - 2) & \text{if }x+2t \ge 2.
        \end{cases}
    \end{equation*}
    This map is continuous by the pasting lemma,
    as $f(x+2t)$, $h(x+2t-1)$, and $g(x+2t-2)$
    are continuous maps on
    $\{(x,t)\in I\times I : x+2t\le 1\}$,
    $\{(x,t)\in I\times I : 1\le x+2t\le 2\}$,
    and $\{(x,t)\in I\times I : x+2t\ge 2\}$, respectively,
    and the maps agree when these sets intersect because
    $f(1) = h(0) = x_1$ and $h(1) = g(0) = y_0$.
    Furthermore $F(x,0) = f(x)$ and $F(x,1) = g(x)$ for all $x\in I$,
    so $F$ is a homotopy from $f$ to $g$.
    Thus any $f,g\colon I\to Y$ are homotopic,
    so $[I,Y]$ has only one element.
\end{soln}

\begin{problem}
    [Munkres 51.3]
    A space $X$ is \vocab{contractible} if the
    identity map $i_X\colon X\to X$ is nulhomotopic.
    \begin{enumerate}[(a)]
        \ii Show that $I$ and $\RR$ are contractible.
	\ii Show that a contractible space is path connected.
    \end{enumerate}
\end{problem}

\begin{soln}
    (a) Let $X = I$ or $\RR$ and consider the map
    $F\colon X\times I\to X$ given by $F(x,t) = (1-t)x$.
    Then $F(x,0) = (1-0)x = x$ and $F(x,1) = (1-1)x = 0$ for all $x\in X$,
    and $F$ is continuous because it is the product of continuous maps
    $(x,t)\mapsto x$ and $(x,t)\mapsto t$,
    so $F$ is a homotopy between the constant map and the zero map
    (this is the same logic as part (a) of the previous problem).
    \\

    (b) Suppose $i_X$ is homotopic to the constant map $f$
    that is equal to $c\in X$ on all of $X$.
    Then, for any $x_1,x_2\in X$, define
    $f_1\colon I\to X$ and $f_2\colon I\to X$ by
    \begin{equation*}
        f_1(t) = F(x_1,t),\quad f_2(t) = F(x_2,1-t)
    \end{equation*}
    so $f_1$ is a path from $x_1$ to $c$
    and $f_2$ is a path from $c$ to $x_2$;
    then $f_1\ast f_2$ is a path from $x_1$ to $x_2$.
\end{soln}

\begin{problem}
    [Munkres 52.3]
    Let $x_0$ and $x_1$ be points of the path-connected space $X$.
    Show that $\pi_1(X,x_0)$ is abelian if and only if
    for every pair $\alpha$ and $\beta$ of paths
    from $x_0$ to $x_1$, we have $\widehat\alpha = \widehat\beta$.
\end{problem}

\begin{soln}
    By Munkres's Theorem 52.1, $\widehat \alpha$ and $\widehat \beta$
    are isomorphisms $\pi_1(X,x_0)\to\pi_1(X,x_1)$,
    so $\widehat\beta\inv\circ\widehat\alpha\colon\pi_1(X,x_0)\to\pi_1(X,x_0)$,
    which is equal to the map
    \begin{equation*}
	[f]\mapsto [\beta * \ol\alpha][f][\alpha * \ol\beta]
    \end{equation*}
    (note that the right-hand side is a product of three elements of $\pi_1(X,x_0)$).
    If $\pi_1(X,x_0)$ is abelian, then this composed map
    sends $[f]\in\pi_1(X,x_0)$ to $[\beta*\ol\alpha][f][\alpha*\ol\beta]
    = [\beta*\ol\alpha][\alpha*\ol\beta][f] = [f]$,
    i.e. the inverse of $\widehat\beta$ is equal to the inverse of $\widehat\alpha$,
    which implies $\widehat\alpha = \widehat\beta$.

    On the other hand suppose $\widehat\alpha = \widehat\beta$
    for any two paths $\alpha,\beta$ from $x_0$ to $x_1$.
    Let $[f],[g]\in\pi_1(X,x_0)$ and $\beta$ an arbitrary path
    from $x_0$ to $x_1$. Define $\alpha = g*\beta$,
    which is also a path from $x_0$ to $x_1$,
    so that $[\alpha*\ol\beta] = [g]$
    (hence also $[\beta*\ol\alpha] = [\ol g]$).
    Then, by assumption, $\widehat\alpha = \widehat\beta$,
    so the two maps share an inverse, meaning
    \begin{equation*}
	[f] = \widehat\beta\inv\widehat\alpha([f])
	= [\beta*\ol\alpha][f][\alpha*\ol\beta]
	= [\ol g][f][g]
	\iff [g][f] = [f][g],
    \end{equation*}
    hence $\pi_1(X,x_0)$ is abelian.
\end{soln}

\begin{problem}
    [Munkres 52.4]
    Let $A\subseteq X$; suppose $r\colon X\to A$ is
    a continuous map such that $r(a) = a$ for each $a\in A$.
    If $a_0\in A$, show that
    \begin{equation*}
        r_*\colon\pi_1(X,a_0)\to\pi_1(A,a_0)
    \end{equation*}
    is surjective.
\end{problem}

\begin{soln}
    For each path $f$ in $A$ based at $a_0$,
    the path homotopy class $[f]$ in $\pi_1(X,a_0)$ gets sent by $r_*$ to
    $r_*([f]) = [r\circ f] = [f]\in\pi_1(A,a_0)$,
    where $r\circ f = f$ because $r$ is the identity on $A$,
    so each $[f]\in\pi_1(A,a_0)$ is mapped to by the corresponding
    path homotopy class $[f]$ in $\in\pi_1(X,a_0)$.
\end{soln}

\begin{problem}
    [Munkres 52.5]
    Let $A$ be a subspace of $\RR^n$
    and let $h\colon(A,a_0)\to(Y,y_0)$ be a basepoint-preserving
    continuous map. Show that if $h$ is extendable to
    a continuous map of $\RR^n$ into $Y$,
    then $h_*$ is the trivial homomorphism.
\end{problem}

\begin{soln}
    Let $\widetilde h\colon(\RR^n,a_0)\to(Y,y_0)$ be a continuous extension of $h$
    and let $i\colon(A,a_0)\to(\RR^n,a_0)$ be the inclusion map.
    Then $\tilde h\circ i = h$, so $h_* = \widetilde h_*\circ i_*$.
    Since $\pi_1(\RR^n,a_0)$ is the trivial group
    (as stated in Section 52 Example 1),
    $\widetilde h_*\colon\pi_1(\RR^n,a_0)\to\pi_1(Y,y_0)$
    is the trivial homomorphism, hence so is $h_* = \widetilde h_*\circ i_*$.
\end{soln}

\begin{problem}
    [Munkres 52.6]
    Let $X$ be path connected. Let $h\colon X\to Y$ be continuous,
    with $h(x_0) = y_0$ and $h(x_1) = y_1$.
    Let $\alpha$ be a path in $X$ from $x_0$ to $x_1$,
    and let $\beta = h\circ\alpha$. Show that
    \begin{equation*}
	\widehat\beta\circ(h_{x_0})_* = (h_{x_1})_*\circ\widehat\alpha.
    \end{equation*}
\end{problem}

\begin{soln}
    Let $[f]\in\pi_1(X,x_0)$. Then,
    \begin{equation*}
	(\widehat\beta\circ(h_{x_0})_*)([f]) = [\ol\beta][h\circ f][\beta]
	= [h\circ\ol\alpha][h\circ f][h\circ\alpha]
	= (h_{x_1})_*([\ol\alpha][f][\alpha])
	= ((h_{x_1})_*\circ\widehat\alpha)([f]).\qedhere
    \end{equation*}
\end{soln}










\end{document}

